% Appendix F -- the staging/zero-copy crossover in full, moved out of section 8
% (2026-09-26) for the page limit. The body states the ordering and why each
% system is right for its own design point; this is the argument behind it.

\section{Staging versus zero-copy, in full}
\label{app:coherence}

Our host tier stages: a layer's KV is copied from a pinned host mirror into a
device rolling slot before attention reads it. DirectKV's kernel instead reads
CPU-resident KV directly, with no staging buffer. We built on a staged backend
for two practical reasons --- it runs on PCIe as well as on a coherent link,
which is what let us measure the negative regime at all, and their kernel is
SM90a-only with no build on our development hardware. Our results are therefore
\emph{conservative with respect to mechanism}: a better transfer path would
raise them.

It is still tempting to read the buffer as a tax that zero-copy removes, and
that reading is wrong in a way that decides which mechanism belongs where. The
two designs hide the transfer behind different things. A zero-copy kernel
overlaps its fetch with the arithmetic of the attention tile it is already
computing --- a window of one tile, and decode attention's intensity is the GQA
group size, 4--8 FLOP/byte against a balance point near 100, so there is almost
no arithmetic there to hide behind. Staging overlaps the fetch for layer
$\ell{+}1$ with the entire compute of layer $\ell$, including the QKV, output
and MLP projections, which are weight-bound and long --- 21 of a 34\,ms step for
Qwen2.5-32B by the projection of \S\ref{sec:model}. \textbf{The staging buffer
buys a layer-sized hiding window instead of a tile-sized one.}

What it costs is two extra HBM passes per byte and the slot's own footprint, so
the comparison is a crossover in offload ratio rather than a dominance. When the
per-step transfer fits inside the layer compute window, staging hides it and pays
only the extra HBM traffic, which at $\Xlink/\BW \approx 0.105$ is a small
fraction of something already hidden; when the transfer exceeds that window it is
exposed under either design, and zero-copy wins by avoiding the passes. This is
why the two systems make opposite choices and both are right: DirectKV places
100\% of KV on the host, so its per-step transfer is the whole working set and can
never fit a compute window, while we place a selected subset and size it to fit
by construction (\S\ref{sec:design:admission}). \textbf{The offload ratio is the
quantity our policy controls}, so the choice of mechanism follows from the policy
rather than constraining it.

A coherent link would help one cost that is not about hiding. Each step projects
new $K,V$ for every host-tier request and we copy them device-to-host to keep the
mirror authoritative; mapped host memory would let the kernel write them in
place, as DirectKV's Alg.~3 does. The saving is small in time --- an asynchronous
version measured 1.6\,ms --- but it removes a class of state rather than an
amount of time, and that state is where migration bugs live. What coherence does
not buy is a cheaper transport: it is a guarantee about visibility, and the same
bytes still cross the link at link bandwidth when the kernel touches them.



\subsection*{Generality, and placement by hardware}

\subsection{Generality across memory tiers and attention backends}

Nothing in the design is specific to HBM-plus-DRAM or to one attention kernel.
The system needs three things from a substrate: a memory larger and slower than
the device's, a path to it whose bandwidth can be measured, and an attention
implementation that can read a block table. A third tier --- CXL-attached
memory, or NVMe --- adds a second budget and a second affinity threshold and
leaves the structure unchanged. What does not survive is the hiding argument: a
tier whose transfer cannot fit inside a layer's compute window is a capacity
tier for idle requests, not a bandwidth tier for active ones, and
Eq.~\ref{eq:gate} decides which it is. Similarly, a zero-copy kernel can replace
our staging buffer through the same backend interface without touching the
policy; the two are ordered by operating point rather than by quality
(\S\ref{sec:discussion:coherence}).

Moving to a new platform requires re-measuring five quantities: achieved
interconnect bandwidth, achieved device bandwidth for weights and for paged KV,
and the host path's two costs, the fixed $F$ and the per-request $c$. The last
two are what the second gate turns on, and $F$ is the one that decides it: the
margins in Table~\ref{tab:gain} are $0.5$ and $2.8$\,ms of $F$, so a platform
whose launch cost differs from ours by a few milliseconds moves rows across the
line in either direction. The rules do not change; their inputs do. This is
why we report the fixed cost separately from the per-byte cost rather than as a
single slope --- only one of the two shrinks when the interconnect improves, and
a system calibrated on their sum will mispredict itself on the next generation.

\subsection{Staging versus zero-copy}
\label{sec:discussion:coherence}

Our host tier stages KV through a device rolling slot; DirectKV's kernel reads
CPU-resident KV directly. We built on a staged backend because it runs on PCIe
as well as on a coherent link --- which is what let us measure the negative
regime at all --- and because their kernel is SM90a-only. Our results are
therefore \emph{conservative with respect to mechanism}.

The two are ordered by operating point rather than by quality, and the ordering
matters for where each belongs. Staging overlaps a layer's fetch with the whole
of the previous layer's compute; a zero-copy kernel overlaps it only with the
attention tile it is already computing, which at decode's arithmetic intensity
is almost no compute at all. Staging buys a layer-sized hiding window and pays
two extra HBM passes for it, so it wins while the per-step transfer fits inside
a layer's compute and loses once the transfer exceeds it. That is why DirectKV,
which places 100\% of KV on the host and therefore can never fit a compute
window, is right to use zero-copy, and why we --- sizing a selected subset to fit
by construction --- are right to stage. Appendix~\ref{app:coherence} works the
crossover through and notes the one cost a coherent link would remove
regardless: the device-to-host write of each step's new $K,V$.

\subsection{Placement by hardware, policy by scheduler}

A further consequence, which we flag as future work because we cannot test it.
With a single coherent copy, demoting a request to the host tier becomes a
policy decision rather than a data movement, and promotion could be lazy and
partial --- the hardware's page migration relocating hot blocks into device
memory while the memory system supplies hints rather than issuing transfers.
That is an appealing division of labour, and to our knowledge nobody has
evaluated it for KV cache.

We are wary of it for one specific reason. Hardware page migration operates at
4\,KB or 2\,MB granularity and migrates on access, which is structurally the
same reactive, fine-grained policy whose migration overhead dominates in the
token-granularity placement literature~\cite{calkv}. Our own claim is that
request-granularity movement avoids that thrash because whole requests move
once, at controlled moments. Handing placement back to a fault-driven mechanism
gives that property up, and whether the hardware's migration is cheap enough to
survive it is an empirical question we have not answered.

% Moved out of section 8 (2026-09-27) when the body hit its page limit. The
% limitations themselves remain in the body; this is the detail behind them.
\subsection*{Generality, and what a new tier must supply}
\label{sec:discussion:general}

Nothing in this paper is specific to host DRAM or to one attention kernel. A tier
must be larger than device memory, expose a path whose bandwidth we can measure at
runtime, and carry a request's per-step traffic inside one layer's compute --- the
last decides whether it can hold \emph{decoding} requests or only parked ones. A
third tier below the second adds a budget and a threshold and leaves the structure
intact. Moving platform means re-measuring five quantities: the two read bandwidths,
the link, and the host path's fixed and per-request costs. We measure those on PCIe
and model NVLink-C2C; whether the argument carries to CXL-attached memory or another
vendor's coherent link depends on the same five, and we have measured none of them
there.

\subsection*{The limitations at length}

\paragraph{Why one GPU.} Tensor parallelism would introduce collectives into
precisely the metric under measurement --- the per-step time whose growth the whole
argument is about --- so a TP baseline would confound the quantity we report with
the collective's cost. Adding a second GPU is also the most obvious
\emph{alternative} to a second memory tier rather than a configuration of it, which
is why we cannot price one against the other on a single-device measurement.

\paragraph{Why the model does not cover MoE or MLA.} Eq.~\ref{eq:step} charges every
weight once per step. For a mixture-of-experts model only the routed experts are
read, so the weight term is a function of the routing distribution rather than the
checkpoint size, and both $\mathrm{ITL}_0$ and $\HBM/\Ckv$ change meaning. An
MLA-class model~\cite{deepseekv2} compresses KV by roughly an order of magnitude,
which raises $B_{\mathrm{cap}}$ enough that the fixed cost is amortised over many
more requests while the ratio the gain depends on is unchanged --- a different
regime, not a scaled one. Neither follows from a reparameterisation of what we
derive, and we have not done the rederivation.

\paragraph{Uniform context, and host DRAM bandwidth.} The cancellation of
\S\ref{sec:model:cancel} is exact only in the fluid limit; both cohorts are counts
of whole requests, so at long context the integer effects the ratio ignores become
the leading correction. And the model charges nothing for host DRAM bandwidth, which
on a 480\,GB Grace configuration is of the same order as the link itself and would
bound $\Xlink$ from behind --- so a measured $\Xlink$ on such a part is an upper
bound on what a loaded host would deliver.

\paragraph{Two engineering limits.} Speculative decoding is unsupported on the host
path, which assumes one query token per step: a rolling slot is sized and rotated
against a single new $K,V$ per layer, and accepting $n$ draft tokens per step would
need $n$ of them. And the pinned host mirrors are real memory, sized to the configured
maximum concurrency and context rather than to live demand, so a node must budget for
them whether or not the tier is in use.

\paragraph{What the alternatives cost.} An operator facing a configuration that
clears both gates can quantize the weights or add a GPU, and both are common. We
think the regime is nonetheless worth serving: quantization is an accuracy decision
a serving team often cannot take unilaterally, and a second GPU doubles the cost of
the node to recover memory that is already installed and idle beside it.
